% Copyright (c) 2026 Geovana Neves. Licensed under CC BY 4.0 (https://creativecommons.org/licenses/by/4.0/). See LICENSE-AND-AUTHORSHIP.md.
%
% 01-workspaces/text/03-inputs.tex

\section{The input files}

\begin{frame}{The row names things, the files define them}
\slidefig[0.62]{guide-inputs-into-run}
\begin{takeaway}By name, never by path: \code{REF}, \code{SET} and \code{PPROC}
carry a file stem, \code{GEOMETRY} names the mesh. That is what lets one matrix
drive two workspaces.\end{takeaway}
\end{frame}

\begin{frame}{\code{matriz.fs}: nineteen columns}
% \relax first: a brace group right after the title is read as a subtitle.
\relax
{\ftstablesize
\begin{tabular}{@{}ll@{}}
\toprule
\textbf{Group} & \textbf{Columns} \\
\midrule
identity: does it run at all & \code{POL}, \code{HIDDEN}, \code{RUN} \\
what it is     & \code{AIRCRAFT}, \code{CONFIGURATION}, \code{DESCRIPTION} \\
what it flies  & \code{FLIGHT\_CONDITION}, \code{SWEEP\_VALUES} \\
what it cites  & \code{GEOMETRY}, \code{REF}, \code{SET}, \code{PPROC} \\
how it solves  & \code{SYMMETRY}, \code{SYMMETRY\_LOADS}, \code{NCPUS}, \code{WALLTIME}, \code{FS\_BUILD} \\
what builds it & \code{WORKFLOW}, \code{VAR\_NAMES\_VALUES} \\
\bottomrule
\end{tabular}}
\vspace{3mm}
\begin{itemize}
  \item A key \textbf{every row} must answer is a column; a key only some rows
        need goes in the free cell, \code{VAR\_NAMES\_VALUES}.
  \item \code{-} means \emph{stated nothing}, so it cannot be confused with a
        truncated line.
\end{itemize}
\end{frame}

\begin{frame}[fragile]{One POL, once, in the whole workspace}
\begin{itemize}
  \item \code{POL} names \code{sims/sim\_<POL>} and the run ids, so it appears
        \textbf{once across every matrix} of the workspace, \code{RUN = 0}
        rows included.
  \item \pyfs{plan} refuses a repeated POL and lists where it appears.
\end{itemize}
\begin{lstlisting}[style=ftsbase]
pyfs-matrix plan matriz.fs --workspace . --fs-version 26.123 --update-ids
\end{lstlisting}
\begin{itemize}
  \item \code{--update-ids} gives each repeated row of \emph{this} matrix the
        next free POL above every POL, run record and \code{sims/} folder
        \textbf{it can see}.
\end{itemize}
\begin{alertblock}{It writes the matrix before planning}
The renumbering stays written even if the plan then refuses. Commit first, so
the change is one diff you can read.
\end{alertblock}
\end{frame}

\begin{frame}[fragile]{The free cell: move, turn, raw}
\begin{lstlisting}[style=ftsbase]
TRANSLATE: {DISTANCE: 0.05 / AXIS: PUSHER_SMRP-X / ALIAS: PUSHER}
ROTATE:    {ANGLE: 3 / AXIS: PUSHER_SMRP-Y / ALIAS: PUSHER}
RAW:       {COMMAND: SOLVER_SET_ITERATIONS 350 / BEFORE: init}
\end{lstlisting}
\begin{columns}[T]
\begin{column}{0.48\textwidth}
\begin{itemize}\small
  \item \code{TRANSLATE}: metres, along \textbf{one} axis; a diagonal
        is two records.
  \item \code{ROTATE}: degrees, about one axis.
  \item \textbf{Every translation comes before every rotation}, whatever the
        order written.
  \item One row is one position: neither is swept.
\end{itemize}
\end{column}
\begin{column}{0.48\textwidth}
\begin{alertblock}{What to know before using it}
\begin{itemize}\small
  \item The alias's frames move with it; \code{<ALIAS>\_SMRP\_ORIGINAL} keeps
        the hub where it was.
  \item A moved part \textbf{comes away} from what it touches: the junction
        is open.
  \item The distance is read in metres; keep the meshes in metres.
\end{itemize}
\end{alertblock}
\end{column}
\end{columns}
\end{frame}

\begin{frame}[fragile]{The flight condition: constraints on one flow state}
\begin{columns}[T]
\begin{column}{0.50\textwidth}
\begin{lstlisting}[style=ftsbase]
FLIGHT_CONDITION: MACH:0.20, REmi:5.5
FLIGHT_CONDITION: TASmps:68.08, ALTFT:10000, dISA:5
\end{lstlisting}
{\ftstablesize
\begin{tabular}{@{}>{\raggedright\arraybackslash}p{0.16\linewidth}>{\raggedright\arraybackslash}p{0.24\linewidth}>{\raggedright\arraybackslash}p{0.46\linewidth}@{}}
\toprule
\textbf{Key} & \textbf{Unit} & \textbf{Constrains} \\
\midrule
\code{MACH}   & none & velocity, by the state's speed of sound \\
\code{TASmps} & m/s & velocity, directly \\
\code{REmi}   & \textbf{millions} & density \\
\code{ALTFT}  & \textbf{feet} & pressure and temperature \\
\code{dISA}   & \textbf{Celsius, a delta} & temperature \\
\bottomrule
\end{tabular}}
\end{column}
\begin{column}{0.46\textwidth}
\begin{itemize}\small
  \item The keys decide what is solved for: the first line solves the
        \textbf{density}, the second is an atmosphere point whose Reynolds
        number follows [5, 7, 8].
  \item Exactly one of \code{MACH} or \code{TASmps}. \code{REmi} needs the
        row's \code{REF}: a Reynolds number has no meaning without a length.
  \item Five keys \textbf{pin} the fluid (\code{RHOkgm3}, \code{MUPas},
        \code{ASMPS}, \code{TK}, \code{PPA}), on the row or in the setup's
        \code{[flight\_condition]}.
  \item A rotor: \(V = J\,(\mathrm{RPM}/60)\,D\); state two of the three.
  \item The set is closed: an unknown key is refused, never ignored.
\end{itemize}
\end{column}
\end{columns}
\end{frame}

\begin{frame}[fragile]{A matrix from an older layout}
\begin{lstlisting}[style=ftsbase]
pyfs-matrix upgrade matriz.fs              # to standard output: read the diff
pyfs-matrix upgrade matriz.fs --in-place   # then write it
\end{lstlisting}
\begin{itemize}\small
  \item A matrix written under an earlier layout is recognised by its header
        and \textbf{refused rather than misread}; \code{upgrade} converts it.
        It needs no recipe, no build and no executable.
  \item A sweep code becomes the word \code{sweep} where the value would be:
        \code{AL} becomes \code{ALPHA:sweep}. A row that varies both angles is
        refused by name: it is one row per sideslip, each with its own POL.
  \item \textbf{The values move across verbatim; the result may not.} A
        Reynolds number that was only recorded becomes a constraint that
        solves for density, so the upgraded row emits a fluid state it never
        emitted before. To keep the old behaviour, drop \code{REmi} and let the
        Reynolds number follow.
\end{itemize}
\end{frame}

\begin{frame}[fragile]{\code{r001.toml}: what things are called (Guide 3)}
\begin{columns}[T]
\begin{column}{0.46\textwidth}
\begin{lstlisting}[style=ftsbase]
area_m2 = 19.502
chord_m =  1.322
span_m  = 15.059
[moment_point]
x_m = 0.0

[aliases]
airframe = ["Spinner","Nacelle"]

[PUSHER]
kind = "rotor"
axis = "X"
diameter_m = 1.80
blade1 = { azimuth_deg = 0, zero = "Y" }
families_blades = ["Blade1", "Blade2"]
\end{lstlisting}
{\tiny\color{SoftGray} The template's synthetic numbers; a study replaces every one.}
\end{column}
\begin{column}{0.50\textwidth}
\begin{itemize}\small
  \item Reference lengths and the moment point divide every coefficient.
  \item An \textbf{alias} is a word rows, groups and plots cite.
  \item A rotor block creates \code{<NAME>\_SMRP}, \code{<NAME>\_RMRP} and
        \code{<NAME>\_RMRP<k>}.
  \item The blade count is the length of \code{families\_blades}.
\end{itemize}
\begin{alertblock}{The failure nothing refuses}
A wrong moment point gives plausible, wrong coefficients. Get it right by
reading.
\end{alertblock}
\end{column}
\end{columns}
\end{frame}

\begin{frame}[fragile]{Setup, post-processing and profiles}
\begin{columns}[T]
\begin{column}{0.32\textwidth}
\begin{block}{\code{s001.toml}}
\small How hard to solve and when to stop. Multiplatform. A command every row
needs goes in its \code{[[raw]]} table.
\end{block}
\end{column}
\begin{column}{0.32\textwidth}
\begin{block}{\code{p001.toml}}
\small What the run exports. The log is declared \textbf{last}: its presence
says the run finished. The \code{.fsm} of every point is always saved;
\code{simulation = false} is refused.
\end{block}
\end{column}
\begin{column}{0.32\textwidth}
\begin{block}{\code{inputs/profiles/}}
\small Probe surveys a \code{[[probes]]} entry cites by \code{points\_file}.
Resolved to an absolute path; a missing file is refused at plan.
\end{block}
\end{column}
\end{columns}
\begin{lstlisting}[style=ftsbase]
[[probes]]
frame = "MRP"
parameters = ["VELOCITY"]
points_file = "disk_survey.txt"     # inputs/profiles/disk_survey.txt
\end{lstlisting}
\end{frame}

\begin{frame}[fragile]{The solver path and the cluster profile}
\begin{columns}[T]
\begin{column}{0.48\textwidth}
\begin{block}{Two files, one sentence}
\small \code{executables.toml} is about the \textbf{study} and is versioned.
\code{executables.local.toml} is about the \textbf{machine} and is not.
\end{block}
\begin{alertblock}{Exactly one HPC profile}
\small More than one in \code{inputs/hpc/} is refused: there is no selector.
\end{alertblock}
\end{column}
\begin{column}{0.48\textwidth}
\begin{lstlisting}[style=ftsbase]
# inputs/hpc/h001.toml
[descriptor.fields]
version = "{fs_build_alias}"

[builds]
"26.123" = "26.1"
\end{lstlisting}
\begin{itemize}\small
  \item The row keeps the build, the profile gives the scheduler's word.
  \item A declaration, not a check: the build number in the collected log is
        the proof.
\end{itemize}
\end{column}
\end{columns}
\end{frame}

\begin{frame}[fragile]{The row states its window}
\begin{itemize}
\item \code{unsteady\_rotor}: \code{LAST\_REVS\_AVG}, in revolutions;
fractions are accepted.
\item \code{unsteady}: \code{LAST\_ITERS\_AVG}, in solver steps.
\item With no window key, \pyfs{plan} refuses the row. \code{WINDOW\_*}
keys are refused, naming the key to write.
\end{itemize}
\small Temporal fragments for the free cell, beside the other required keys:
\begin{lstlisting}[style=ftsbase]
DELTA_THETA: 5 / REVOLUTIONS: 4 / LAST_REVS_AVG: 1
DELTA_TIME: 0.01 / TIME_ITERATIONS: 100 / LAST_ITERS_AVG: 20
\end{lstlisting}
\begin{alertblock}{One declaration of each attitude}
\small \code{ALPHA} or \code{BETA} in both \code{FLIGHT\_CONDITION} and
\code{VAR\_NAMES\_VALUES} is refused, even if the values agree.
\end{alertblock}
\end{frame}

\begin{frame}[fragile]{Groups and optional pproc tables}
\begin{columns}[T]
\begin{column}{0.47\textwidth}
\begin{lstlisting}[style=ftsbase]
[groups]
PUSHER = "PUSHER"
AIRFRAME = "airframe"
TOTAL = "all"

[phase_locked]
min_revolutions = 4.0
last_revolutions_avg = 2.0

[products]
superfile_format = "legacy_polar"
\end{lstlisting}
\end{column}
\begin{column}{0.49\textwidth}
\begin{itemize}\small
\item One alias as a string. A list is refused with the
one-alias line to write; several members go in the reference's \code{[aliases]}.
\item \code{[phase\_locked]}: one row per azimuth across revolutions;
no declaration keeps the passage series.
\item \code{[equations]} and \code{[glossary]} bind too: derived
columns at post, named \code{<NAME>\_<alias>}; equations may chain.
\item Fixed-width super file: sixteen characters per field; \code{csv}
remains the default.
\end{itemize}
\end{column}
\end{columns}
\end{frame}

\begin{frame}[fragile]{Names and plot-group declarations}
\begin{lstlisting}[style=ftsbase]
[names]
CL_MRP_TOTAL = "CL_TOTAL"
\end{lstlisting}
\begin{itemize}\small
\item Requires that plotted source column. The whole dictionary applies or
none does; missing sources and occupied destinations produce a \code{\#names}
skip. \code{CL} is already the native last-step health check.
\item Applied after axes and equations to the unsteady polar, time average
and passage series. Plots, per-blade and azimuthal headings retain their names.
\item Plot-group names: letters, digits, underscores and bare
\code{\{family\}} only. No spaces, format spec or conversion.
\item Ambiguous names cannot supply rotor histories or global axis blocks.
\item \code{VARIABLES.md} and \code{WRITING-EQUATIONS.md} are generated in
\code{inputs/pproc/} by init, plan and post; commit or ignore them.
\end{itemize}
\end{frame}

\begin{frame}[fragile]{A mesh file, OBJ or STL}
\begin{columns}[T]
\begin{column}{0.50\textwidth}
\begin{lstlisting}[style=ftsbase]
# wing/wing.boundaries.toml, beside wing.obj
boundaries = ["Wing"]      # OBJ: from its groups

[import]
units = "METER"            # the FILE's unit

[trailing_edges]
file = "wing.te.txt"       # points file
\end{lstlisting}
\begin{alertblock}{Export the mesh in metres}
\small On 26.124 a wing OBJ in millimetres, its trailing edge detected,
solved as the same wing in metres to the print [3]: \code{IMPORT}
converts. The file route in millimetres, and every other build, are not
measured.
\end{alertblock}
\end{column}
\begin{column}{0.46\textwidth}
\begin{itemize}\small
  \item \code{GEOMETRY: wing.obj}, like a \code{.fsm}. The script starts a
        new simulation, imports in the stated unit, then works in metres.
  \item No \code{[import]}: refused at plan. An OBJ with no sidecar gets
        one from the plan, its \code{boundaries} read from the file's groups.
        \textbf{Still by hand}: the \code{[import]} unit, the
        trailing edge, and an STL's names, where a name in the wrong place
        cites the wrong surface.
  \item \code{[[import.operations]]}: scale, rename, mirror, translate,
        rotate, in the order written.
\end{itemize}
\end{column}
\end{columns}
\end{frame}

\begin{frame}[fragile]{The trailing edge, by file}
\begin{lstlisting}[style=ftsbase]
METER                                   # the unit of the points
0.216872113,-0.002159825,0.223199974    # one MID-POINT of an edge per line
\end{lstlisting}
\begin{itemize}\small
  \item \textbf{The default route.} Every point is checked against the mesh at
        plan; the first one off an edge is refused by line and position.
  \item The run writes the node file in metres, imports it with
        \code{IMPORT\_WAKE\_EDGES\_FROM\_FILE}, and compares the solver's count
        with the points: a difference is \code{FAILED\_SCRIPT}. A point that
        misses an edge is otherwise \textbf{silent}.
  \item \textbf{26.124 only} [3]. The row must export its log.
  \item Detection is the second route, written in full: \code{detect = "auto"},
        or by surface with a sweep angle.
  \item \code{BASE\_REGIONS} names the \textbf{base}, not the body [3].
\end{itemize}
\end{frame}

\begin{frame}{A public airframe to practise on}
\begin{columns}[T]
\begin{column}{0.54\textwidth}
\begin{itemize}\small
  \item The SMI airframe geometry [9] is published as IGES components: the
        fuselage, the nacelles with their pylons, the lifting surfaces, and
        the wing-body geometry they share. Each component is whole, so
        intermediate configurations can be built from it.
  \item It carries \textbf{no propeller}. Any rotor is added by the user and
        placed by a rotor block of the reference (Guide 3), from any
        propeller geometry the user holds.
  \item Mesh it with the site's mesher, export OBJ \textbf{in metres}, and
        write its \code{<stem>.boundaries.toml} beside it: the route of the
        two previous frames.
\end{itemize}
\end{column}
\begin{column}{0.42\textwidth}
\begin{block}{Its licence}
\small CC BY-NC 4.0: credit the record, no commercial use. Cite it as [9]
wherever a result uses it.
\end{block}
\begin{takeaway}A first workspace on a known public geometry makes a
mistake visible before it costs a campaign.\end{takeaway}
\end{column}
\end{columns}
\end{frame}

\begin{frame}[fragile]{Every key in one generated page, and an optional workbook}
\begin{columns}[T]
\begin{column}{0.48\textwidth}
\begin{block}{\code{inputs/INPUTS.md} and \code{inputs/input\_template.md}}
\begin{itemize}\small
  \item \code{INPUTS.md} lists every key an input may state: the row by run
        type, the setup, the pproc, the reference, the sidecar. Each says what
        it sets, its unit or values and the command it reaches.
  \item \code{input\_template.md} shows every input \textbf{file}, one
        complete example per kind, each one a file the package reads as it
        stands.
  \item Both are \textbf{generated} by \pyfs{init}, \pyfs{plan} and
        \pyfs{post}; never edit them.
\end{itemize}
\end{block}
\end{column}
\begin{column}{0.48\textwidth}
\begin{block}{An Excel workbook, if you want one [6]}
\small An ordinary \code{.xlsx}, no macros. Python creates it and
synchronises saved files in both directions, always as a preview first, then
an apply:
\begin{lstlisting}[style=ftsbase,basicstyle=\ttfamily\tiny]
pip install "pyflightstream[excel]"
python -m pyflightstream.workspace.excel preview
   <book>.xlsx --workspace . --direction read
   --batch read-preview.json
python -m pyflightstream.workspace.excel apply
   read-preview.json
\end{lstlisting}
The ASCII matrix stays the record; the workbook is an editor.
\end{block}
\end{column}
\end{columns}
\end{frame}

\section{Site folders}

\begin{frame}{\code{ansa/}, \code{builds/}, \code{docs/}}
\begin{columns}[T]
\begin{column}{0.32\textwidth}
\begin{block}{\code{ansa/}}
\small The mesh scripts. The script is versioned, the mesh is not.
\end{block}
\end{column}
\begin{column}{0.32\textwidth}
\begin{block}{\code{builds/}}
\small One folder per solver build number.
\end{block}
\end{column}
\begin{column}{0.32\textwidth}
\begin{block}{\code{docs/}}
\small What a second person needs to repeat the first one's work: these
guides, the install notes.
\end{block}
\end{column}
\end{columns}
\vspace{2mm}
\begin{block}{Adding a build, in order}
\small \textbf{1} copy into its own folder \quad \textbf{2} register in
\code{executables.toml} \quad \textbf{3} point your local file at it, and add its
\code{[builds]} line on the cluster \quad \textbf{4} plan one row and read the
build number in its log. \textbf{Keep the old build.}
\end{block}
\end{frame}
